% GENERATED FILE. Edit canonical lessons, then run npm run generate:books.
% canonical-source-hash: fnv1a64:4648aa62e6ebcf60
% canonical-lessons: PT-C290-arma, PT-C290-pistola, PT-C290-ameacar, PT-C290-ameaca, PT-C290-assassino

\chapter{Armas e ameaças}
\label{ch:pt-armas-e-ameacas}
\hlchaptermodality{290}

\begin{chapteropening}
\textbf{By the end of this chapter:} \emph{I can talk about weapons and threats, and name a killer.}

Complete the last lesson of chapter 290: I can talk about weapons and threats, and name a killer.
\end{chapteropening}

\section[a arma]{a arma — the weapon, the gun}

\label{lesson:PT-C290-arma}



\begin{quote}
Before the new one: say the Portuguese for the con man, then the Portuguese for the accomplice.
\end{quote}

\subsection*{You'll want to know: a arma}
\textbf{a arma} — \textquotedblleft{}the weapon, the gun\textquotedblright{}.

\begin{quote}
A polícia encontrou a arma no rio.
\end{quote}

\begin{quote}
\emph{The police found the weapon in the river.}
\end{quote}

\emph{Uma arma de fogo} — a firearm.

\begin{cousinweb}
\emph{arma} is inherited from Latin \emph{arma} \textquotedblleft{}tools, arms\textquotedblright{} (EtymDB), a plural that Portuguese came to treat as a singular.
\end{cousinweb}

\begin{grammarlens}[title={What you've built}]
One more word for telling, explaining and describing at B1.
\end{grammarlens}

\subsection*{Your turn}
\begin{itemize}
\raggedright
  \item \emph{Say it:} \emph{a arma}
  \item \emph{Say it:} \emph{a arma}, once more
  \item \emph{Say it:} a sentence of your own with \emph{a arma}
  \item \emph{Recall:} say the Portuguese for the con man, then the Portuguese for the accomplice, then say \emph{a arma} again
\end{itemize}

\begin{tcolorbox}[breakable,colback=teal!4,colframe=teal!35!black,title={Before you move on}]
What does a arma mean? (\textquotedblleft{}The weapon, the gun\textquotedblright{}.) Say it once more.
\end{tcolorbox}
\section[a pistola]{a pistola — the pistol, the handgun}

\label{lesson:PT-C290-pistola}



\begin{quote}
Before the new one: say the Portuguese for the accomplice, then the Portuguese for the weapon.
\end{quote}

\subsection*{You'll want to know: a pistola}
\textbf{a pistola} — \textquotedblleft{}the pistol, the handgun\textquotedblright{}.

\begin{quote}
O assaltante tinha uma pistola.
\end{quote}

\begin{quote}
\emph{The robber had a pistol.}
\end{quote}

Also a spray gun: \emph{uma pistola de pintura}.

\begin{cousinweb}
\emph{pistola} was borrowed through German \emph{Pistole} (EtymDB), from Czech \emph{píšťala} \textquotedblleft{}a pipe, a whistle\textquotedblright{}: the long thin barrel looked like a pipe.
\end{cousinweb}

\begin{grammarlens}[title={What you've built}]
One more word for telling, explaining and describing at B1.
\end{grammarlens}

\subsection*{Your turn}
\begin{itemize}
\raggedright
  \item \emph{Say it:} \emph{a pistola}
  \item \emph{Say it:} \emph{a pistola}, once more
  \item \emph{Say it:} a sentence of your own with \emph{a pistola}
  \item \emph{Recall:} say the Portuguese for the accomplice, then the Portuguese for the weapon, then say \emph{a pistola} again
\end{itemize}

\begin{tcolorbox}[breakable,colback=teal!4,colframe=teal!35!black,title={Before you move on}]
What does a pistola mean? (\textquotedblleft{}The pistol, the handgun\textquotedblright{}.) Say it once more.
\end{tcolorbox}
\section[ameaçar]{ameaçar — to threaten}

\label{lesson:PT-C290-ameacar}



\begin{quote}
Before the new one: say the Portuguese for the weapon, then the Portuguese for the pistol.
\end{quote}

\subsection*{You'll want to know: ameaçar}
\textbf{ameaçar} — \textquotedblleft{}to threaten\textquotedblright{}.

\begin{quote}
O homem ameaçou o empregado com uma faca.
\end{quote}

\begin{quote}
\emph{The man threatened the employee with a knife.}
\end{quote}

\emph{Ameaçar de morte} — to threaten to kill.

\begin{cousinweb}
\emph{ameaçar} comes down from Old Portuguese \emph{amẽaçar} (EtymDB), built on Vulgar Latin \emph{minācia} \textquotedblleft{}a threat\textquotedblright{}, from Latin \emph{mināx} \textquotedblleft{}threatening\textquotedblright{}.
\end{cousinweb}

\begin{grammarlens}[title={What you've built}]
One more word for telling, explaining and describing at B1.
\end{grammarlens}

\subsection*{Your turn}
\begin{itemize}
\raggedright
  \item \emph{Say it:} \emph{ameaçar}
  \item \emph{Say it:} \emph{ameaçar}, once more
  \item \emph{Say it:} a sentence of your own with \emph{ameaçar}
  \item \emph{Recall:} say the Portuguese for the weapon, then the Portuguese for the pistol, then say \emph{ameaçar} again
\end{itemize}

\begin{tcolorbox}[breakable,colback=teal!4,colframe=teal!35!black,title={Before you move on}]
What does ameaçar mean? (\textquotedblleft{}To threaten\textquotedblright{}.) Say it once more.
\end{tcolorbox}
\section[a ameaça]{a ameaça — the threat}

\label{lesson:PT-C290-ameaca}



\begin{quote}
Before the new one: say the Portuguese for the pistol, then the Portuguese for to threaten.
\end{quote}

\subsection*{You'll want to know: a ameaça}
\textbf{a ameaça} — \textquotedblleft{}the threat\textquotedblright{}.

\begin{quote}
Ela recebeu uma ameaça por telefone.
\end{quote}

\begin{quote}
\emph{She received a threat by phone.}
\end{quote}

\emph{Sob ameaça} — under threat.

\begin{cousinweb}
\emph{ameaça} comes from Vulgar Latin \emph{minācia} \textquotedblleft{}a threat\textquotedblright{} (EtymDB), from Latin \emph{mināx} \textquotedblleft{}threatening\textquotedblright{}.
\end{cousinweb}

\begin{grammarlens}[title={What you've built}]
One more word for telling, explaining and describing at B1.
\end{grammarlens}

\subsection*{Your turn}
\begin{itemize}
\raggedright
  \item \emph{Say it:} \emph{a ameaça}
  \item \emph{Say it:} \emph{a ameaça}, once more
  \item \emph{Say it:} a sentence of your own with \emph{a ameaça}
  \item \emph{Recall:} say the Portuguese for the pistol, then the Portuguese for to threaten, then say \emph{a ameaça} again
\end{itemize}

\begin{tcolorbox}[breakable,colback=teal!4,colframe=teal!35!black,title={Before you move on}]
What does a ameaça mean? (\textquotedblleft{}The threat\textquotedblright{}.) Say it once more.
\end{tcolorbox}
\section[o assassino]{o assassino — the murderer, the killer}

\label{lesson:PT-C290-assassino}



\begin{quote}
Before the new one: say the Portuguese for to threaten, then the Portuguese for the threat.
\end{quote}

\subsection*{You'll want to know: o assassino}
\textbf{o assassino} — \textquotedblleft{}the murderer, the killer\textquotedblright{}.

\begin{quote}
Nos filmes, o assassino acaba por voltar ao mesmo lugar.
\end{quote}

\begin{quote}
\emph{In films, the killer ends up going back to the same place.}
\end{quote}

A woman is \emph{a assassina}.

\begin{cousinweb}
\emph{assassino} came through Italian \emph{assassino}, which is usually traced to Arabic \emph{ḥaššāšīn}, the name of a medieval sect; EtymDB has the Portuguese word with no source.
\end{cousinweb}

\begin{grammarlens}[title={What you've built}]
One more word for telling, explaining and describing at B1.
\end{grammarlens}

\subsection*{Your turn}
\begin{itemize}
\raggedright
  \item \emph{Say it:} \emph{o assassino}
  \item \emph{Say it:} \emph{o assassino}, once more
  \item \emph{Say it:} a sentence of your own with \emph{o assassino}
  \item \emph{Recall:} say the Portuguese for to threaten, then the Portuguese for the threat, then say \emph{o assassino} again
\end{itemize}

\begin{tcolorbox}[breakable,colback=teal!4,colframe=teal!35!black,title={Before you move on}]
What does o assassino mean? (\textquotedblleft{}The murderer, the killer\textquotedblright{}.) Say it once more.
\end{tcolorbox}
